% ============================================================
% Section 4: Directional Relational Manifold (DRM)
% ============================================================

\section{Directional Relational Manifold (DRM)}
\label{sec:drm}

The Directional Relational Manifold (DRM) is the geometric substrate
on which the epistemic system operates. Each token is mapped to a point
on a learned 5D Riemannian manifold, where geodesic distances encode
epistemic relationships.

\paragraph{Relation to information geometry.}
Information geometry \cite{Amari2000} endows the space of probability
distributions with a Riemannian structure via the Fisher Information Matrix
$g_{ij}(\theta) = \E[\partial_i \log p \cdot \partial_j \log p]$.
This framework, intrinsic by Chentsov's theorem \cite{Chentsov1982}, unifies
statistical estimation and divergence geometry under a single formalism
\cite{Nielsen2020}.

The DRM used in \aletheion{} extends information geometry to the
variable-dimension setting: where Amari's framework operates on manifolds of
fixed dimension~$n$, DRM allows $\dim\mathcal{D}(p)$ to vary
pointwise---modelling the epistemic head's ability to represent varying
numbers of active uncertainty dimensions depending on the input.  When
$\dim\mathcal{D}(p) = n$ (constant) and $g_p$ is the Fisher metric, DRM
reduces to standard information geometry.

\subsection{Manifold Embedding}

The mapping from hidden states to 5D coordinates is performed by an
MLP with sigmoid activation on the last layer:

\begin{equation}
    \bm{c} = \sigma\big(\bm{W}_2 \cdot \text{GELU}(\bm{W}_1 \bh + \bm{b}_1) + \bm{b}_2\big) \in [0, 1]^5
    \label{eq:manifold_emb}
\end{equation}

with $\bm{W}_1 \in \R^{h \times d}$, $\bm{W}_2 \in \R^{5 \times h}$,
where $h = 5 \times 4 = 20$ (configurable via \texttt{hidden\_factor}),
totaling {\raise.17ex\hbox{$\scriptstyle\sim$}}16K parameters for $d = 768$.

The embedding also computes L2 distances to 6 fixed
\textbf{anchor points} in the manifold (truth, ignorance, noise,
complex, stale, ideal), refined via a small projection layer.

The 5 dimensions capture distinct epistemic aspects:

\begin{table}[H]
\centering
\caption{Interpretation of the 5 manifold dimensions.}
\label{tab:dims_5d}
\begin{tabular}{cl}
\toprule
\textbf{Dim.} & \textbf{Interpretation} \\
\midrule
$c_0$ & Factual certainty \\
$c_1$ & Logical consistency \\
$c_2$ & Context / grounding \\
$c_3$ & Confidence calibration \\
$c_4$ & Epistemic depth \\
\bottomrule
\end{tabular}
\end{table}

\subsection{Metric Tensor}

The manifold geometry is defined by a metric tensor
$\bG \in \R^{5 \times 5}$ that must be \textbf{symmetric positive-definite}
(SPD). We construct $\bG$ from a learned lower triangular matrix $\bm{L}$:

\begin{equation}
    \bG = \bm{L}\bm{L}^\top + \epsilon \bm{I}_5
    \label{eq:metric_spd}
\end{equation}

where $\bm{L}$ is parametrized by $5 \times (5+1) / 2 = 15$ free
parameters and $\epsilon = 10^{-6}$ ensures numerical stability. The
Cholesky decomposition guarantees that $\bG$ is SPD by construction.

\begin{proposition}[Metric Tensor Positivity]
For any lower triangular $\bm{L} \in \R^{5 \times 5}$ with non-zero
diagonal, $\bG = \bm{L}\bm{L}^\top + \epsilon \bm{I}$ is SPD.
\end{proposition}

\begin{proof}
For any $\bm{v} \neq 0$:
$\bm{v}^\top \bG \bm{v} = \bm{v}^\top \bm{L}\bm{L}^\top \bm{v} + \epsilon \|\bm{v}\|^2
= \|\bm{L}^\top \bm{v}\|^2 + \epsilon \|\bm{v}\|^2 > 0$.
\end{proof}

\subsection{Geodesic Distance}

The geodesic distance between two points $\bm{c}_i, \bm{c}_j$ on the
manifold is approximated by the Mahalanobis distance induced by $\bG$:

\begin{equation}
    d_{\bG}(\bm{c}_i, \bm{c}_j) = \sqrt{(\bm{c}_i - \bm{c}_j)^\top \bG\, (\bm{c}_i - \bm{c}_j)}
    \label{eq:geodesic}
\end{equation}

This is a first-order approximation of the true geodesic distance, valid
when the curvature is locally uniform. In the general case, the geodesic
would require integrating the differential equation:

\begin{equation}
    \ddot{c}^\mu + \Gamma^\mu_{\alpha\beta} \dot{c}^\alpha \dot{c}^\beta = 0
\end{equation}

where $\Gamma^\mu_{\alpha\beta}$ are the Christoffel symbols.

\subsection{Directional Field}

The directional field $\bm{F} \in \R^5$ captures the preferred direction
of evolution on the manifold, derived from the entropy of attention
patterns:

\begin{equation}
    H_\ell = -\sum_t \bm{A}_{\ell,t} \log(\bm{A}_{\ell,t} + \epsilon)
\end{equation}

\begin{equation}
    \bm{F} = \text{MLP}_F(\bar{H}) \in \R^5
\end{equation}

where $\bar{H}$ is the mean entropy across layers and heads.

\subsection{Ricci Curvature}

Manifold health can be assessed by the Ricci scalar curvature, computed
from the metric tensor:

\begin{equation}
    R = g^{\mu\nu} R_{\mu\nu} = g^{\mu\nu} R^\lambda_{\ \mu\lambda\nu}
\end{equation}

In practice, we use a computationally efficient approximation based on
the condition number of $\bG$:

\begin{equation}
    \text{health}_{\text{metric}} = 1 - \frac{\log(\kappa(\bG))}{\log(\kappa_{\max})}
\end{equation}

where $\kappa(\bG) = \lambda_{\max} / \lambda_{\min}$ is the condition
number. A healthy manifold has $\kappa(\bG) \approx 1$ (isotropic
geometry).

\subsection{Metric Tensor Regularization}

To prevent geometric degeneration, we add a regularization on the
condition number:

\begin{equation}
    \calL_{\text{metric}} = \log\big(\kappa(\bG) + 1\big)
\end{equation}

This forces $\bG$ to remain close to the identity, preventing certain
manifold directions from collapsing or inflating excessively.
The logarithmic form provides a softer penalty than $(\kappa - 1)^2$,
allowing the metric to deviate slightly from isotropy while preventing
extreme degeneration.
